\section{Boundary accounting for duon current}\label{sec:boundary-accounting}

Boundary accounting records current and closure across declared structures. It acts on fields, regions, boundaries, shells, windows, and fractal addresses while preserving the ontology of those structures.

\paragraph{Regional boundary template.}
A reusable regional statement has the form
\[
\mathbf a\vdash(\mathcal F,\mathcal R,\partial R,\omega),
\]
where \(\mathcal F\) is the observer-field, \(\mathcal R\) is the declared retained region, \(\partial R\) is the region boundary, and \(\omega\) is the window. Slice flux, window current, and moving-boundary flux use the same objects in every example below.

\begin{definition}[Duon slice flux]
For a declared boundary \(\partial R\), define the duon-current slice flux by
\begin{equation}
\Phi^D_{\partial R}(t)
=
\sum_{e\in\partial R(t)}
\epsilon_e C_{3,e}.
\label{eq:slice-flux}
\end{equation}
Here \(e\) is a duon segment crossing or returning at the boundary, \(\epsilon_e\) is signed orientation, and \(C_{3,e}\) is retained trit-capacity carried by the segment. When a fractal address must be shown, write \(\mathbf a\vdash \Phi^D_{\partial R}(t)\).
\end{definition}

\begin{definition}[Window duon current]
The window current is
\begin{equation}
J^D_{\partial R}(\omega)
=
\sum_{t\in\omega}
\Phi^D_{\partial R}(t).
\label{eq:window-current}
\end{equation}
\end{definition}

\paragraph{Flux/current distinction.}
A slice flux is one declared cut-frame crossing of duon/trit current through the boundary. A window current is the retained many-frame continuation of such crossings across the declared window:
\[
\Phi^D_{\partial R}(t)=\text{one cut-frame crossing},
\]
\[
J^D_{\partial R}(\omega)
=
\sum_{t\in\omega}\Phi^D_{\partial R}(t)
=\text{retained many-frame continuation}.
\]

\paragraph{Throughput fractions.}
Let \(J^D_{\mathrm{abs}}\) be the absolute throughput of duon/trit current on the boundary. Then
\[
\eta_{\mathrm{in}}
=
\frac{J^D_{\mathrm{in}}}{J^D_{\mathrm{abs}}},
\qquad
\eta_{\mathrm{out}}
=
\frac{J^D_{\mathrm{out}}}{J^D_{\mathrm{abs}}},
\qquad
\eta_{\mathrm{pending}}
=
\frac{J^D_{\mathrm{pending}}}{J^D_{\mathrm{abs}}}.
\]

\paragraph{Moving boundary.}
For a moving retained boundary, the window may record crossing, sweep, and shedding terms as
\[
\Phi^D_{\partial R,\mathrm{mov}}(t)
=
\Phi^D_{\mathrm{cross}}(t)
-
\Phi^D_{\mathrm{sweep}}(t)
+
\Phi^D_{\mathrm{shed}}(t).
\]
Here crossing is duon current crossing the retained boundary, sweep is support swept by boundary motion, and shedding is disturbance that breaks from the retained boundary.

\begin{figure}[H]
\centering
\fbox{\begin{minipage}{0.92\linewidth}
\centering
\textbf{Duon-current boundary accounting}\\[0.4em]
\begin{tabular}{c c c c}
\textbf{Internal retained} & \textbf{Outward shedding} & \textbf{Reflected return} & \textbf{Pending return}\\
recurring knot/core & crosses \(\partial R\) & returns through shell/target & remains unclosed over \(\omega\)\\
scalar burden & shell crossing & updates next emission & duration-window overlap\\
\end{tabular}\\[0.8em]
\(\Phi^D_{\partial R}(t)=\sum_{e\in\partial R(t)}\epsilon_e C_{3,e}\),\quad
\(J^D_{\partial R}(\omega)=\sum_{t\in\omega}\Phi^D_{\partial R}(t)\)\\[0.4em]
Scalar side: \(P^D\) = duonic pressure.\qquad Vector side: \(A^{(q)}\) = field attention vector.
\end{minipage}}
\caption{Duon-current boundary accounting. A declared boundary \(\partial R\) separates internal retained current, outward shedding, reflected return, and pending return current.}
\label{fig:v39-duon-current-boundary-accounting}
\end{figure}

\subsection{Boundary-accounting examples}\label{subsec:boundary-accounting-examples}
The same boundary/current form applies across familiar region choices. Declare a region, declare \(\partial R\), and account for current crossing, returning, and remaining pending.

\begin{figure}[H]
\centering
\fbox{\begin{minipage}{0.90\linewidth}
\centering
\textbf{Water--atmosphere boundary accounting}\\[0.4em]
\begin{tabular}{c c c}
Water region & \(\partial R\) interface & Atmosphere\\
inner vortex/core current & in/out/reflected current counted & outward shedding\\
retained pressure & interface duration & returned coupling\\
\end{tabular}
\end{minipage}}
\caption{Water--atmosphere boundary accounting. The interface \(\partial R\) separates lower-region retained current, outward shedding into the upper region, reflected return, and pending duration.}
\label{fig:v39-water-atmosphere-example}
\end{figure}

\begin{figure}[H]
\centering
\fbox{\begin{minipage}{0.90\linewidth}
\centering
\textbf{Wire current and field interaction}\\[0.4em]
\begin{tabular}{c c c}
Conductor core & \(\partial R\) cylindrical shell & Field shell\\
axial retained current & boundary accounting & surrounding return paths\\
core pressure & pending duration shell & reflected coupling\\
\end{tabular}
\end{minipage}}
\caption{Wire-current and field interaction. A conductor region carries core current while the surrounding shell carries reflected return and pending return burden.}
\label{fig:v39-wire-current-example}
\end{figure}

\begin{figure}[H]
\centering
\fbox{\begin{minipage}{0.90\linewidth}
\centering
\textbf{Plasma and magnetic-field coupling}\\[0.4em]
\begin{tabular}{c c c}
Plasma core & \(\partial R\) magnetic shell & Coupled field shell\\
circulating retained current & shell crossing counted & delayed return paths\\
plasma pressure & pending duration & field attention vector\\
\end{tabular}
\end{minipage}}
\caption{Plasma and magnetic-field coupling. A plasma core carries internal duon current; the magnetic shell carries reflected and pending-duration duon-current components across \(\partial R\).}
\label{fig:v39-plasma-example}
\end{figure}
